\section[Deep-learned scatter initialisation]{Deep-Learned Scatter Initialisation for \glsentryshort{sirf}}
\label{app:scatter_dl}

The bremsstrahlung reconstructions of Chapters~\ref{chapterlabel3} and~\ref{chapter:paper} initialise the preparatory additive-term estimation using a convolutional neural network that predicts scatter projections directly from the measured emission projections and \glsentryshort{ct}-derived attenuation information, following Xiang et al. \cite{xiangDeepNeuralNetwork2020}. The code used in this thesis is a fork of the original authors' implementation, adapted to run inference within \gls{sirf} workflows \cite{ovtchinnikovSIRFSynergisticImage2020}: it adds \glsentryshort{sirf}-facing configuration and inference scripts driven by a YAML configuration file, alongside the pre-trained network weights and example notebooks. As discussed in Chapter~\ref{chapterlabel3}, the pre-trained network was developed under a different scanner and protocol configuration than the datasets used in this thesis; the scatter-derived output is therefore used only as a heuristic warm start for the broader non-primary term. It is refined by the damped reconstruct--simulate loop driven by \texttt{simind-python-connector} (Section~\ref{app:simind_python_connector}), and the final estimate is frozen before the target reconstruction.

The adapted repository is available at \href{https://github.com/samdporter/spect-scatter-deep-learning}{GitHub} and retains the structure of the upstream implementation by Xiang et al., to which all credit for the network architecture and training belongs; only the \glsentryshort{sirf} integration and inference path are contributions of this thesis. \TODO{Archive a tagged release of \texttt{spect-scatter-deep-learning} on Zenodo and replace this GitHub link with the DOI-backed citation.}
